\documentclass[11pt,a4paper]{article}
\usepackage[utf8]{inputenc}
\usepackage[T1]{fontenc}
\usepackage[portuguese]{babel}
\usepackage{geometry}
\geometry{margin=2.5cm}
\usepackage{listings}
\usepackage{xcolor}
\usepackage{amsmath}
\usepackage{amssymb}
\usepackage{hyperref}
\usepackage{enumitem}
\usepackage{tcolorbox}
\usepackage{tabularx}
\usepackage{booktabs}
\usepackage{graphicx}

\definecolor{codebg}{rgb}{0.95,0.95,0.95}
\definecolor{codegreen}{rgb}{0.1,0.5,0.1}
\definecolor{codepurple}{rgb}{0.5,0.1,0.5}
\definecolor{codegray}{rgb}{0.5,0.5,0.5}

\lstdefinestyle{python}{
    backgroundcolor=\color{codebg},
    commentstyle=\color{codegray}\itshape,
    keywordstyle=\color{codepurple}\bfseries,
    stringstyle=\color{codegreen},
    basicstyle=\ttfamily\small,
    breaklines=true,
    frame=single,
    language=Python,
    showstringspaces=false,
    tabsize=4,
    literate={á}{{\'a}}1 {ã}{{\~a}}1 {é}{{\'e}}1 {í}{{\'i}}1 {ó}{{\'o}}1 {õ}{{\~o}}1 {ú}{{\'u}}1 {ç}{{\c{c}}}1 {Á}{{\'A}}1 {Ã}{{\~A}}1 {É}{{\'E}}1 {Í}{{\'I}}1 {Ó}{{\'O}}1 {Õ}{{\~O}}1 {Ú}{{\'U}}1 {Ç}{{\c{C}}}1 {â}{{\^a}}1 {ê}{{\^e}}1 {ô}{{\^o}}1 {û}{{\^u}}1 {Â}{{\^A}}1 {Ê}{{\^E}}1 {Ô}{{\^O}}1 {Û}{{\^U}}1 {à}{{\`a}}1 {è}{{\`e}}1 {ì}{{\`i}}1 {ò}{{\`o}}1 {ù}{{\`u}}1 {ä}{{\"a}}1 {ö}{{\"o}}1 {Ä}{{\"A}}1 {Ö}{{\"O}}1
}
\lstset{style=python}

\newtcolorbox{objetivos}[1][]{
  colback=green!5!white,
  colframe=green!40!black,
  title=O que você vai aprender nesta página,
  fonttitle=\bfseries,
  #1
}

\newtcolorbox{exercicio}[1]{
  colback=blue!5!white,
  colframe=blue!40!black,
  title=#1,
  fonttitle=\bfseries
}

\newtcolorbox{solucao}{
  colback=yellow!5!white,
  colframe=orange!60!black,
  title=Solução proposta,
  fonttitle=\bfseries
}

\title{\textbf{Data Analysis with Python 2024--2025}\\Parte 6: Machine Learning Types\\\large Principal Component Analysis (Teoria)}
\author{Tradução e Adaptação: University of Helsinki --- MOOC.fi}
\date{}

\begin{document}

\maketitle

\section*{Nota Importante}
Este material é uma tradução e adaptação baseada no curso \textit{Data Analysis with Python 2024-2025}, oferecido pela \textbf{Universidade de Helsinque (MOOC.fi)}.

\tableofcontents
\newpage

\section{Machine Learning: Análise de Componentes Principais (PCA)}

A análise de componentes principais (PCA) é um método de aprendizado não supervisionado (\emph{unsupervised learning}) que tenta detectar as direções nas quais os vetores formados por um conjunto de dados variam mais.

Ele primeiro encontra a direção de maior variância e, em seguida, prossegue para descobrir novas direções de maior variância que sejam \textbf{ortogonais} àquelas já encontradas. Assim, para dados de $n$ dimensões, ele retorna (por padrão) $n$ direções ortogonais e as variâncias correspondentes a elas. Estas direções são chamadas de \textbf{eixos principais} (\emph{principal axes}), e se projetarmos um ponto do conjunto de dados sobre estes eixos, obtemos os \textbf{componentes principais} (\emph{principal components}) de cada eixo.

Para usar outra terminologia, o conjunto de eixos principais forma uma \emph{base} para o espaço vetorial onde os pontos de dados residem, e os componentes principais representam as coordenadas dos pontos neste novo sistema de coordenadas.

A classe \texttt{PCA} na biblioteca \texttt{scikit-learn} possui o método \texttt{transform}, que converte dados originais para este novo sistema de coordenadas.

\subsection{Exemplo de Distribuição Multivariada}
Vamos analisar um exemplo de distribuição Gaussiana multivariada submetida ao processo de PCA. Primeiro, nós extraímos aleatoriamente dados de uma distribuição, espalhamos os pontos e visualizamos através de \texttt{matplotlib}.

\begin{lstlisting}
import numpy as np
import matplotlib.pyplot as plt
import math
from sklearn.decomposition import PCA

rng = np.random.RandomState(0)
X = rng.randn(2, 400)
scale = np.array([[1, 0], [0, 0.4]]) # Desvios Padroes: 1 e 0.4
rotate = np.array([[1, -1], [1, 1]]) / math.sqrt(2)

transform = np.dot(rotate, scale)
X = np.dot(transform, X).T

# O PCA revelara os eixos principais desse amontoado
pca = PCA(n_components=2)
pca.fit(X)

print("Eixos principais:\n", pca.components_)
print("Variancia Explicada:\n", pca.explained_variance_)
\end{lstlisting}

\subsubsection{Redução de Dimensionalidade}

Você pode ter notado que podemos fornecer ao construtor PCA um parâmetro \texttt{n\_components}. Esse parâmetro fornece o número de eixos que queremos manter. Se o valor for $n$, o algoritmo retorna os $n$ componentes com as mais altas variâncias e ignora ativamente os demais. Dessa forma, esse algoritmo também age como uma exímia técnica de \textbf{redução de dimensionalidade}.

\begin{lstlisting}
pca = PCA(n_components=1)
pca.fit(X)
Z = pca.transform(X) # Transforma os pontos para caber em uma dimensao
\end{lstlisting}

A redução pode ser usada, por exemplo, para converter dados hiperdimensionais abstratos para 2 ou 3 dimensões, garantindo que o humano seja capaz de visualizá-los em um gráfico compreensível. Além disso, essa técnica atua perfeitamente como \textbf{pré-processamento}: você pode rodar o PCA antes de injetar os dados em Regressões Lineares e Redes Neurais para reduzir o ruído (fatores que variam pouco) e preservar o cerne dos atributos do conjunto de treinamento.

\section{Resumo da Semana 6}

\begin{itemize}
    \item Conhecemos um método de aprendizado supervisionado: a classificação \textbf{Naive Bayes}. Usamos \textit{cross validation} para garantir que o modelo não está viciado (\emph{overfitting}).
    \item Na sessão de \emph{Clustering} (Agrupamento), vimos os exemplos não-supervisionados de \textbf{K-Means}, \textbf{DBSCAN}, e agrupamento hierárquico, todos baseados na noção de distância entre nós de dados.
    \item E, por fim, a análise de \textbf{PCA} demonstrou uma forma não-supervisionada eficiente para tratar reduções dimensionais filtrando matrizes pela sua capacidade variável explicativa.
\end{itemize}

\end{document}
